% Input-process-output model -- include with % Input-process-output model -- include with % Input-process-output model -- include with % Input-process-output model -- include with \input{chapters/ipo_model}
% Needs only TikZ + the positioning library (already loaded via forest in main.tex)
\begin{figure}[H]
  \centering
  \begin{tikzpicture}[
      >=latex,
      font=\small,
      io/.style  ={draw, rounded corners=4pt, minimum width=2.6cm, minimum height=1.1cm},
      proc/.style={draw, fill=blue!6, minimum width=2.8cm, minimum height=1.1cm},
      note/.style={font=\scriptsize\itshape, text=black!60, align=center},
      flow/.style={->, thick},
    ]
    \node[io]                    (in)   {Input};
    \node[proc, right=1.6cm of in]  (proc) {Process};
    \node[io,   right=1.6cm of proc] (out)  {Output};

    % Side bars: flowchart symbol for a (predefined) process
    \draw ([xshift=4pt]proc.north west) -- ([xshift=4pt]proc.south west);
    \draw ([xshift=-4pt]proc.north east) -- ([xshift=-4pt]proc.south east);

    \draw[flow] (in)   -- (proc);
    \draw[flow] (proc) -- (out);

    % Example sources and sinks
    \node[note, below=3pt of in]  {e.g.\ sensors,\\sub-systems};
    \node[note, below=3pt of out] {e.g.\ actuators,\\sub-systems};
  \end{tikzpicture}

  \caption[Input--process--output model]{%
    Input--process--output model. The process receives its input from
    other systems such as sensors and passes its output on to further
    sub-systems such as actuators. As the only element between input and
    output, the process is a single point of failure (SPOF).}
  \label{fig:ipo-model}
\end{figure}

% Needs only TikZ + the positioning library (already loaded via forest in main.tex)
\begin{figure}[H]
  \centering
  \begin{tikzpicture}[
      >=latex,
      font=\small,
      io/.style  ={draw, rounded corners=4pt, minimum width=2.6cm, minimum height=1.1cm},
      proc/.style={draw, fill=blue!6, minimum width=2.8cm, minimum height=1.1cm},
      note/.style={font=\scriptsize\itshape, text=black!60, align=center},
      flow/.style={->, thick},
    ]
    \node[io]                    (in)   {Input};
    \node[proc, right=1.6cm of in]  (proc) {Process};
    \node[io,   right=1.6cm of proc] (out)  {Output};

    % Side bars: flowchart symbol for a (predefined) process
    \draw ([xshift=4pt]proc.north west) -- ([xshift=4pt]proc.south west);
    \draw ([xshift=-4pt]proc.north east) -- ([xshift=-4pt]proc.south east);

    \draw[flow] (in)   -- (proc);
    \draw[flow] (proc) -- (out);

    % Example sources and sinks
    \node[note, below=3pt of in]  {e.g.\ sensors,\\sub-systems};
    \node[note, below=3pt of out] {e.g.\ actuators,\\sub-systems};
  \end{tikzpicture}

  \caption[Input--process--output model]{%
    Input--process--output model. The process receives its input from
    other systems such as sensors and passes its output on to further
    sub-systems such as actuators. As the only element between input and
    output, the process is a single point of failure (SPOF).}
  \label{fig:ipo-model}
\end{figure}

% Needs only TikZ + the positioning library (already loaded via forest in main.tex)
\begin{figure}[H]
  \centering
  \begin{tikzpicture}[
      >=latex,
      font=\small,
      io/.style  ={draw, rounded corners=4pt, minimum width=2.6cm, minimum height=1.1cm},
      proc/.style={draw, fill=blue!6, minimum width=2.8cm, minimum height=1.1cm},
      note/.style={font=\scriptsize\itshape, text=black!60, align=center},
      flow/.style={->, thick},
    ]
    \node[io]                    (in)   {Input};
    \node[proc, right=1.6cm of in]  (proc) {Process};
    \node[io,   right=1.6cm of proc] (out)  {Output};

    % Side bars: flowchart symbol for a (predefined) process
    \draw ([xshift=4pt]proc.north west) -- ([xshift=4pt]proc.south west);
    \draw ([xshift=-4pt]proc.north east) -- ([xshift=-4pt]proc.south east);

    \draw[flow] (in)   -- (proc);
    \draw[flow] (proc) -- (out);

    % Example sources and sinks
    \node[note, below=3pt of in]  {e.g.\ sensors,\\sub-systems};
    \node[note, below=3pt of out] {e.g.\ actuators,\\sub-systems};
  \end{tikzpicture}

  \caption[Input--process--output model]{%
    Input--process--output model. The process receives its input from
    other systems such as sensors and passes its output on to further
    sub-systems such as actuators. As the only element between input and
    output, the process is a single point of failure (SPOF).}
  \label{fig:ipo-model}
\end{figure}

% Needs only TikZ + the positioning library (already loaded via forest in main.tex)
\begin{figure}[H]
  \centering
  \begin{tikzpicture}[
      >=latex,
      font=\small,
      io/.style  ={draw, rounded corners=4pt, minimum width=2.6cm, minimum height=1.1cm},
      proc/.style={draw, fill=blue!6, minimum width=2.8cm, minimum height=1.1cm},
      note/.style={font=\scriptsize\itshape, text=black!60, align=center},
      flow/.style={->, thick},
    ]
    \node[io]                    (in)   {Input};
    \node[proc, right=1.6cm of in]  (proc) {Process};
    \node[io,   right=1.6cm of proc] (out)  {Output};

    % Side bars: flowchart symbol for a (predefined) process
    \draw ([xshift=4pt]proc.north west) -- ([xshift=4pt]proc.south west);
    \draw ([xshift=-4pt]proc.north east) -- ([xshift=-4pt]proc.south east);

    \draw[flow] (in)   -- (proc);
    \draw[flow] (proc) -- (out);

    % Example sources and sinks
    \node[note, below=3pt of in]  {e.g.\ sensors,\\sub-systems};
    \node[note, below=3pt of out] {e.g.\ actuators,\\sub-systems};
  \end{tikzpicture}

  \caption[Input--process--output model]{%
    Input--process--output model. The process receives its input from
    other systems such as sensors and passes its output on to further
    sub-systems such as actuators. As the only element between input and
    output, the process is a single point of failure (SPOF).}
  \label{fig:ipo-model}
\end{figure}
